\textbf{Phần 1. Trắc nghiệm lựa chọn}
\setcounterpageref{ex}{0}

\begin{ex} %[3D2N09]
Kết quả của phép nhân $6 \times 4$ là
\choice
{$18$}
{$28$}
{$30$}
{\True $24$}
\loigiai{
$6 \times 4 = 24$.\\
\textbf{Đáp án của câu hỏi là D}
}
\end{ex}

\begin{ex} %[3D2N09]
Kết quả của phép chia $30 : 6$ là
\choice
{$4$}
{\True $5$}
{$6$}
{$7$}
\loigiai{
$30 : 6 = 5$.\\
\textbf{Đáp án của câu hỏi là B}
}
\end{ex}

\begin{ex} %[3D2H09]
Kết quả của phép nhân $6 \times 7$ là
\choice
{$36$}
{$48$}
{$40$}
{\True $42$}
\loigiai{
$6 \times 7 = 42$.\\
\textbf{Đáp án của câu hỏi là D}
}
\end{ex}

\begin{ex} %[3D2H09]
Kết quả của phép chia $48 : 6$ là
\choice
{$7$}
{\True $8$}
{$9$}
{$6$}
\loigiai{
$48 : 6 = 8$.\\
\textbf{Đáp án của câu hỏi là B}
}
\end{ex}

\begin{ex} %[3D2H09]
Điền số thích hợp vào chỗ chấm: $6 \times \dots = 54$.
\choice
{\True $9$}
{$8$}
{$7$}
{$6$}
\loigiai{
$54 : 6 = 9$.\\
\textbf{Đáp án của câu hỏi là A}
}
\end{ex}

\begin{ex} %[3D2H09]
Điền số thích hợp vào chỗ chấm: $\dots : 6 = 6$.
\choice
{$30$}
{$42$}
{\True $36$}
{$24$}
\loigiai{
$6 \times 6 = 36$.\\
\textbf{Đáp án của câu hỏi là C}
}
\end{ex}

\begin{ex} %[3D2V09]
Tính giá trị biểu thức: $6 \times 8 - 18$.
\choice
{\True $30$}
{$28$}
{$32$}
{$36$}
\loigiai{
$6 \times 8 - 18 = 48 - 18 = 30$.\\
\textbf{Đáp án của câu hỏi là A}
}
\end{ex}

\begin{ex} %[3D2V09]
Tính giá trị biểu thức: $42 : 6 + 25$.
\choice
{$30$}
{\True $32$}
{$34$}
{$36$}
\loigiai{
$42 : 6 + 25 = 7 + 25 = 32$.\\
\textbf{Đáp án của câu hỏi là B}
}
\end{ex}

\begin{ex} %[3D2V09]
Một tuần lễ Lan đi học 6 ngày. Hỏi trong 4 tuần lễ như thế Lan đi học bao nhiêu ngày?
\choice
{\True $24$ ngày}
{$20$ ngày}
{$28$ ngày}
{$18$ ngày}
\loigiai{
$6 \times 4 = 24$ ngày.\\
\textbf{Đáp án của câu hỏi là A}
}
\end{ex}

\begin{ex} %[3D2V09]
Có 36 cái bánh xếp đều vào 6 hộp. Mỗi hộp có số cái bánh là
\choice
{$5$ cái}
{$7$ cái}
{\True $6$ cái}
{$8$ cái}
\loigiai{
$36 : 6 = 6$ cái bánh.\\
\textbf{Đáp án của câu hỏi là C}
}
\end{ex}

\begin{ex} %[3D2V09]
Một cuộn vải dài $54\text{ m}$, người ta cắt thành các khúc vải dài $6\text{ m}$. Số khúc vải cắt được là
\choice
{$8$ khúc}
{$7$ khúc}
{$10$ khúc}
{\True $9$ khúc}
\loigiai{
$54 : 6 = 9$ khúc vải.\\
\textbf{Đáp án của câu hỏi là D}
}
\end{ex}

\begin{ex} %[3D2C09]
Có 6 túi kẹo, mỗi túi có 8 cái kẹo. Người ta đã lấy ra 15 cái kẹo để chia cho các bạn. Số kẹo còn lại là
\choice
{$31$ cái}
{$35$ cái}
{\True $33$ cái}
{$29$ cái}
\loigiai{
Tổng số kẹo: $6 \times 8 = 48$ cái.\\
Còn lại: $48 - 15 = 33$ cái kẹo.\\
\textbf{Đáp án của câu hỏi là C}
}
\end{ex}

\textbf{Phần 2. Trắc nghiệm đúng sai}
\setcounterpageref{ex}{0}

\begin{ex} %[3D2H09]
Xét tính đúng sai của các phép tính trong bảng nhân 6, bảng chia 6:
\choiceTF
{\True $6 \times 3 = 18$.}
{\True $24 : 6 = 4$.}
{$6 \times 9 = 56$.}
{\True $6 \times 9 = 54$.}
\loigiai{
\begin{itemize}
    \item $6 \times 3 = 18$: Đúng.
    \item $24 : 6 = 4$: Đúng.
    \item $6 \times 9 = 54 \neq 56$. Ý 3 Sai, Ý 4 Đúng.
\end{itemize}
\begin{itemize}
    \item a) \textbf{Đúng}.
    \item b) \textbf{Đúng}.
    \item c) \textbf{Sai}.
    \item d) \textbf{Đúng}.
\end{itemize}
}
\end{ex}

\begin{ex} %[3D2V09]
Một lớp có 30 học sinh chia đều thành các tổ, mỗi tổ có 6 học sinh. Xét tính đúng sai:
\choiceTF
{\True Lớp đó chia được thành 5 tổ.}
{\True Mỗi tổ có 6 học sinh.}
{Hai tổ có tất cả 10 học sinh.}
{\True Hai tổ có tất cả 12 học sinh.}
\loigiai{
\begin{itemize}
    \item $30 : 6 = 5$ tổ: Đúng.
    \item Mỗi tổ 6 bạn: Đúng.
    \item 2 tổ: $6 \times 2 = 12$ bạn $\neq 10$. Ý 3 Sai, Ý 4 Đúng.
\end{itemize}
\begin{itemize}
    \item a) \textbf{Đúng}.
    \item b) \textbf{Đúng}.
    \item c) \textbf{Sai}.
    \item d) \textbf{Đúng}.
\end{itemize}
}
\end{ex}

\textbf{Phần 3. Trả lời ngắn}
\setcounterpageref{ex}{0}

\begin{ex} %[3D2H09]
Tính tích: $6 \times 5$.
\shortans{30}
\loigiai{
\textbf{Đáp số:} $30$\\
$6 \times 5 = 30$.
}
\end{ex}

\begin{ex} %[3D2H09]
Tính thương: $18 : 6$.
\shortans{3}
\loigiai{
\textbf{Đáp số:} $3$\\
$18 : 6 = 3$.
}
\end{ex}

\begin{ex} %[3D2V09]
Tính giá trị biểu thức: $6 \times 6 + 14$.
\shortans{50}
\loigiai{
\textbf{Đáp số:} $50$\\
$36 + 14 = 50$.
}
\end{ex}

\begin{ex} %[3D2V09]
Một can đựng $6\text{ l}$ mật ong. Hỏi 7 can như thế đựng bao nhiêu lít mật ong?
\shortans{42}
\loigiai{
\textbf{Đáp số:} $42$\\
$6 \times 7 = 42\text{ l}$.
}
\end{ex}

\begin{ex} %[3D2V09]
Có 48 quả cam xếp đều vào 6 đĩa. Hỏi mỗi đĩa có bao nhiêu quả cam?
\shortans{8}
\loigiai{
\textbf{Đáp số:} $8$\\
$48 : 6 = 8$ quả cam.
}
\end{ex}

\begin{ex} %[3D2C09]
Tìm $x$ biết: $x \times 6 = 60 - 18$.
\shortans{7}
\loigiai{
\textbf{Đáp số:} $7$\\
$x \times 6 = 42 \Rightarrow x = 42 : 6 = 7$.
}
\end{ex}
